\documentclass[10pt,a4paper]{article}

% ---------- Packages ----------
\usepackage[utf8]{inputenc}
\usepackage[T1]{fontenc}
\usepackage[english]{babel}

\usepackage[
    a4paper,
    margin=1.8cm,
    top=1.2cm,
    bottom=1.5cm,
    includehead,
    headheight=1.9cm,
    headsep=0.45cm
]{geometry}

\usepackage{graphicx}
\usepackage[table]{xcolor}
\usepackage{titlesec}
\usepackage{enumitem}
\usepackage{fancyhdr}
\usepackage{array}
\usepackage{tabularx}
\usepackage{xurl}
\usepackage{hyperref}

% ---------- Colors ----------
\definecolor{mainblue}{RGB}{25,55,110}
\definecolor{lightgray}{RGB}{245,245,245}

% ---------- Links ----------
\hypersetup{
    colorlinks=true,
    linkcolor=mainblue,
    urlcolor=mainblue,
    citecolor=mainblue
}

% ---------- Header ----------
\pagestyle{fancy}
\fancyhf{}
\renewcommand{\headrulewidth}{0.4pt}
\renewcommand{\footrulewidth}{0pt}

\fancyhead[L]{%
    \begin{minipage}[c]{0.32\textwidth}
        \raggedright
        \IfFileExists{uzi_logo.png}{\includegraphics[width=3.2cm]{\detokenize{uzi_logo.png}}}{\textbf{UZH}}
    \end{minipage}
}

\fancyhead[C]{%
    \begin{minipage}[c]{0.36\textwidth}
        \centering
        \textbf{\small MSc Thesis Proposal}\\[-0.1em]
        \scriptsize Department of Informatics
    \end{minipage}
}

\fancyhead[R]{}

\fancyfoot[C]{\thepage}

% ---------- Section Formatting ----------
\titleformat{\section}
  {\color{mainblue}\normalfont\large\bfseries}
  {}
  {0em}
  {}

\titlespacing*{\section}{0pt}{0.6em}{0.25em}

% ---------- Compact Lists ----------
\setlist[itemize]{noitemsep, topsep=2pt, leftmargin=1.2em}
\setlist[enumerate]{noitemsep, topsep=2pt, leftmargin=1.4em}

% ---------- Tables ----------
\renewcommand{\arraystretch}{1.08}
\setlength{\tabcolsep}{6pt}

\begin{document}

% ---------- Title ----------
\begin{center}
    {\Large \textbf{Evidence-Grounded Fault Diagnosis from Sensor Events and\\[0.15em]Maintenance Documents Using Retrieval-Augmented Generation}}
\end{center}

\vspace{0.35em}

% ---------- Student Information ----------
\noindent
\begin{tabularx}{\textwidth}{>{\bfseries}l X >{\bfseries}l X}
Student Name: & Necati Furkan Colak       & Student ID:     & 24746323 \\
Program:      & MSc Computer Science      & Semester:       & Spring 2026 \\
Supervisor:   & [Supervisor Name]         & Co-Supervisor:  & [Co-Supervisor Name] \\
Email:        & necatifurkan.colak@uzh.ch & Date:           & \today \\
\end{tabularx}

\vspace{0.5em}
\hrule
\vspace{0.5em}

\section{Abstract}

Predictive-maintenance systems output anomaly scores, sensor alarms, and failure predictions. Turning those numbers into likely causes, comparable past failures, and recommended inspection actions still depends on expert reading. Retrieval-augmented generation (RAG) \cite{lewis2020rag} can connect a sensor event to maintenance knowledge and produce an explanation that cites the evidence it used. This thesis builds a RAG system that takes sensor events and anomaly descriptions, retrieves matching fault evidence and maintenance documents, and generates a cited fault explanation. It evaluates how different retrieval strategies affect diagnosis accuracy and how retrieval errors propagate to the final answer. The work uses one open predictive-maintenance dataset whose fault types are turned into synthetic event reports and matched with open manuals, fault tables, and controlled maintenance records. Evaluation covers fault-diagnosis accuracy, retrieval recall, citation correctness, faithfulness, and an error-propagation analysis. The contribution is a measured link between retrieval quality and fault-diagnosis quality, and a reproducible pipeline that grounds fault explanations in traceable sources.

\section{Background and Motivation}

A maintenance system can flag that a machine is likely to fail, but the flag alone does not tell a technician what is wrong or what to check. That step, from a numeric signal to a probable cause and an inspection action, is where most of the expert time goes. RAG can support it by retrieving the fault descriptions, tables, and procedures that match the current sensor event, and by presenting an explanation with citations the technician can verify. The value of such a system depends on retrieving the right evidence, which makes the retrieval step the part worth studying.

\section{Research Gap}

Combining language models with retrieval over event logs, sensor data, and manuals for fault diagnosis has been shown to be workable, but existing accounts bundle many data sources and mix retrieval with separate knowledge-base construction. This leaves the retrieval step itself underexamined. In particular, how retrieval errors carry through to a wrong or unsupported diagnosis is not isolated and measured. A narrower study can focus on sensor-event-to-evidence retrieval and quantify that error propagation \cite{gao2024survey, es2024ragas}.

\section{Research Objective}

Build a RAG system that retrieves relevant fault evidence and maintenance documents from sensor events and anomaly descriptions, and measure how retrieval strategy affects fault-diagnosis accuracy and citation correctness.

\section{Research Questions}

\begin{enumerate}
    \item \textbf{RQ1 (primary):} How effectively can RAG ground fault-diagnosis explanations in sensor events and maintenance documentation?
    \item \textbf{RQ2:} Does hybrid retrieval improve fault-evidence retrieval over dense retrieval, and does sensor-metadata filtering improve diagnosis accuracy and citation correctness?
    \item \textbf{RQ3:} How do retrieval errors propagate to the final diagnosis?
\end{enumerate}

\section{Proposed Methodology}

\begin{itemize}
    \item \textbf{Data:} One open predictive-maintenance dataset, chosen to bound scope, for example the AI4I 2020 dataset \cite{matzka2020} or the NASA C-MAPSS turbofan dataset \cite{saxena2008}. The dataset fault types are converted into synthetic event reports and matched with open manuals, fault tables, and researcher-prepared controlled maintenance records. No proprietary data is required.
    \item \textbf{Architecture:} Sensor-event textualization, fault metadata extraction, hybrid retrieval (BM25 and dense), metadata filtering, reranking, and citation-grounded fault explanation. The contribution is not a new prediction model: ground-truth fault labels from the dataset, or a simple off-the-shelf classifier, supply the fault signal.
    \item \textbf{Baselines:} LLM without retrieval; dense-vector RAG; hybrid and metadata-aware RAG.
    \item \textbf{Evaluation:} Fault-diagnosis accuracy or macro-F1; retrieval Recall@k; evidence coverage; citation correctness; faithfulness \cite{es2024ragas}; unsupported-recommendation rate; an error-propagation analysis that traces retrieval misses to diagnosis failures; latency. The benchmark holds 80 to 150 fault scenarios, split into development and held-out test sets.
\end{itemize}

\section{Expected Contribution}

\begin{enumerate}
    \item \textbf{Academic:} An experimental analysis of the relationship between retrieval quality and fault-diagnosis quality, including how retrieval errors propagate.
    \item \textbf{Technical:} A reproducible pipeline that combines sensor-event data with maintenance-document retrieval.
    \item \textbf{Practical:} A method that produces cited, auditable fault explanations for maintenance personnel.
\end{enumerate}

\section{Scope and Limitations}

The thesis does not build a new remaining-useful-life model, does not optimize maintenance schedules, and does not trigger automatic maintenance actions. It uses one dataset, a limited fault taxonomy, and one RAG prototype. Simulation-model development and multiple machine types are out of scope, which keeps the work narrower than a full fault-diagnosis knowledge-base build.

\section{Feasibility}

Synthetic event reports and a single open dataset keep data preparation and compute within a several-month master period. Retrieval and generation use standard components, and the fault signal comes from existing labels rather than a new model, so the engineering effort centers on the retrieval pipeline and its evaluation.

\section{Preliminary Work Plan}

\begin{tabularx}{\textwidth}{|p{1.9cm}|X|}
\hline
\rowcolor{lightgray}
\textbf{Period} & \textbf{Planned Work} \\
\hline
Month 1 & Dataset and fault taxonomy selection. \\
\hline
Month 2 & Technical document corpus construction. \\
\hline
Month 3 & Sensor-event preprocessing and textualization. \\
\hline
Month 4 & Baselines and proposed RAG implementation. \\
\hline
Month 5 & Experiments and error-propagation analysis. \\
\hline
Month 6 & Thesis writing; final presentation. \\
\hline
\end{tabularx}

\begingroup
\renewcommand{\refname}{\texorpdfstring{\color{mainblue}}{}Preliminary References}
\scriptsize
\begin{thebibliography}{9}
\setlength{\itemsep}{0pt}

\bibitem{lewis2020rag}
P.\ Lewis et al.
\textit{Retrieval-Augmented Generation for Knowledge-Intensive NLP Tasks}.
NeurIPS, 2020.

\bibitem{gao2024survey}
Y.\ Gao et al.
\textit{Retrieval-Augmented Generation for Large Language Models: A Survey}.
arXiv:2312.10997, 2023.

\bibitem{es2024ragas}
S.\ Es et al.
\textit{RAGAS: Automated Evaluation of Retrieval Augmented Generation}.
EACL, 2024.

\bibitem{matzka2020}
S.\ Matzka.
\textit{Explainable Artificial Intelligence for Predictive Maintenance Applications}.
Third International Conference on Artificial Intelligence for Industries (AI4I), 2020.

\bibitem{saxena2008}
A.\ Saxena, K.\ Goebel, D.\ Simon, and N.\ Eklund.
\textit{Damage Propagation Modeling for Aircraft Engine Run-to-Failure Simulation}.
International Conference on Prognostics and Health Management (PHM), 2008.

\bibitem{ragindustry2025}
\textit{Retrieval-Augmented Generation in Industry: An Interview Study on Use Cases, Requirements, Challenges, and Evaluation}.
arXiv:2508.14066, 2025.

\end{thebibliography}
\endgroup

\vspace{0.1em}

\noindent
\begin{tabularx}{\textwidth}{X X}
\textbf{Student Signature:} & \textbf{Supervisor Signature:} \\[0.9em]
\rule{0.42\textwidth}{0.4pt} & \rule{0.42\textwidth}{0.4pt} \\
\end{tabularx}

\end{document}
